\documentclass[../ICIT-Paper-2024.tex]{subfiles}
\begin{document}
\paragraph{}

After collecting the data, we preprocess these data. The purpose of these steps is to identify and adjust outlier values. Handling these outlier values ensures that they do not affect the results of models. The method applied for processing is Winsorization~\cite{dixon1974trimming}, which includes two steps: (1) detect outliers, and (2) calculate the threshold values and transform the data.

\paragraph{Detect Outlier}

We used the percentile method and observed the results to establish threshold limits for the features in the wallets. If the value of a feature exceeds this threshold limit, it will be considered an outlier.

\paragraph{Calculate the cut-off values and transform the data}

After determining the threshold values, Winsorization calculates the values at these percentiles from our dataset. Specifically, values higher than the upper threshold (and lower than the lower threshold) will be changed to the upper threshold value (or the lower threshold value). The table~\ref{tab2} shows the results of determining the upper and lower thresholds of the total assets attribute. Then, the values within the threshold continue to be normalized using min-max normalization to bring them into the range [300, 850] according to the following formula~\ref{eq1}.
\begin{equation}\label{eq1}
X' = \frac{(X - X_{\text{min}})*(850-300)}{X_{\text{max}} - X_{\text{min}}} + 300
\end{equation}

\begin{itemize}
    \item \( X \) is the original value.
    \item \( X' \) is the normalized value.
    \item \( X_{\text{min}} \) is the minimum value of the feature.
    \item \( X_{\text{max}} \) is the maximum value of the feature.
\end{itemize}
\paragraph{Example}
For example, let's consider the processing of a single attribute, which is total assets.
\begin{table}[h!]
\centering
\caption{Percentile threshold and value of total asset attribute}\label{tab2}
\begin{tabular}{|c|c|}
\hline
\textbf{Percentile Threshold} & \textbf{Value} \\
\hline
0 & 0.0 \\
\hline
0.05 & 0.19 \\
\hline
0.25 & 9.61 \\
\hline
0.9 & 4219.90 \\
\hline
0.95 & 11936.94 \\
\hline
1 & 22824864075 \\
\hline
\end{tabular}

\label{tab:example}
\end{table}

The values in the total asset attribute exhibit significant variation. Choosing thresholds of 0.25 and 0.95 corresponds to approximate values of \$10 and \$\numprint{10000}, respectively. These are considered levels to distinguish wallets with no assets from those with a significant amount of assets. Subsequently, standardization is applied.
\end{document}